\documentclass[11pt,a4paper]{article}
\usepackage[utf8]{inputenc}
\usepackage[T1]{fontenc}
\usepackage{lmodern}
\usepackage[french]{babel}
\usepackage{amsmath,amssymb,mathtools}
\usepackage{icomma}
\usepackage{xcolor}
\usepackage[most]{tcolorbox}
\usepackage{enumitem}
\usepackage[a4paper,margin=2.2cm,headheight=26pt,headsep=10pt]{geometry}
\usepackage{fancyhdr}
\usepackage[hidelinks]{hyperref}
\definecolor{primary}{RGB}{222,49,99}
\definecolor{primarydark}{RGB}{150,22,55}
\definecolor{excol}{RGB}{0,135,90}
\tcbset{boxbase/.style={enhanced,breakable,boxrule=0.9pt,arc=2.5pt,
  left=3mm,right=3mm,top=2.3mm,bottom=2.3mm,fonttitle=\bfseries,coltitle=white}}
\newtcolorbox[auto counter]{exo}[1][]{boxbase,colback=primary!4,colframe=primary,
  title={\textsc{Exercice}~\thetcbcounter\ifx\relax#1\relax\relax\else\textmd{ --- \itshape #1}\fi}}
\newcommand{\R}{\mathbb{R}}
\newcommand{\N}{\mathbb{N}}
\newcommand{\Z}{\mathbb{Z}}
\newcommand{\ds}{\displaystyle}
\pagestyle{fancy}\fancyhf{}
\fancyhead[L]{\small\textcolor{primarydark}{\textbf{Thème 1} — Puissances, racines et exposants}}
\fancyhead[R]{\small\textcolor{primarydark}{\textsc{Facile}}}
\fancyfoot[C]{\small\thepage}
\renewcommand{\headrulewidth}{0.8pt}
\begin{document}

\begin{center}
{\LARGE\bfseries\color{primarydark} Niveau Facile — Exercices}\\[2pt]
{\color{primary}\rule{0.45\textwidth}{1.2pt}}\\[3pt]
{\small\itshape Une seule mécanique à la fois. Aucune calculatrice.}
\end{center}
\vspace{1mm}

\begin{exo}[règles produit et quotient, même base]
Écris chaque expression sous la forme d'une \emph{seule} puissance $a^{k}$ (donne l'exposant $k$).
\begin{enumerate}[label=\alph*),leftmargin=2em,topsep=1pt,itemsep=1pt]
  \item $a^{4}\times a^{3}$ \qquad
  \item $\ds\frac{a^{7}}{a^{2}}$ \qquad
  \item $a^{6}\times a^{-2}$ \qquad
  \item $\ds\frac{a^{3}}{a^{8}}$ \qquad
  \item $a^{2}\times a^{5}\times a$
\end{enumerate}
\end{exo}

\begin{exo}[exposant nul et exposant négatif]
Réécris sans exposant négatif (ou calcule la valeur exacte), avec $a\in\R_0$ :
\begin{enumerate}[label=\alph*),leftmargin=2em,topsep=1pt,itemsep=1pt]
  \item $5^{0}$ \qquad
  \item $a^{-3}$ \qquad
  \item $\ds\frac{1}{a^{4}}$ (sous forme $a^{k}$) \qquad
  \item $2^{-3}$ (valeur exacte) \qquad
  \item $\ds\Bigl(\frac{1}{3}\Bigr)^{-2}$ (valeur exacte)
\end{enumerate}
\end{exo}

\begin{exo}[racines et exposants fractionnaires]
Passe de la racine à l'exposant fractionnaire (ou l'inverse), avec $c>0$ :
\begin{enumerate}[label=\alph*),leftmargin=2em,topsep=1pt,itemsep=1pt]
  \item $\sqrt[5]{c}$ \qquad
  \item $\sqrt[3]{c^{2}}$ \qquad
  \item $c^{1/4}$ (en racine) \qquad
  \item $\sqrt[3]{c^{6}}$ (simplifie) \qquad
  \item $\sqrt{c^{10}}$ (simplifie)
\end{enumerate}
\end{exo}

\begin{exo}[signe et parenthésage]
Calcule la valeur exacte, en faisant très attention à la position du signe « $-$ » :
\begin{enumerate}[label=\alph*),leftmargin=2em,topsep=1pt,itemsep=1pt]
  \item $(-3)^{2}$ \qquad
  \item $-3^{2}$ \qquad
  \item $(-2)^{3}$ \qquad
  \item $-2^{4}$ \qquad
  \item $\sqrt{(-5)^{2}}$
\end{enumerate}
\end{exo}

\begin{exo}[équation exponentielle à même base]
Résous dans $\R$ (l'exponentielle est injective : on égale les exposants) :
\begin{enumerate}[label=\alph*),leftmargin=2em,topsep=1pt,itemsep=1pt]
  \item $2^{x}=2^{5}$ \qquad
  \item $3^{2x}=3^{8}$ \qquad
  \item $5^{x+1}=5^{4}$ \qquad
  \item $7^{3x}=7^{0}$ \qquad
  \item $10^{x-2}=10^{3}$
\end{enumerate}
\end{exo}

\end{document}
